\rowcolors{2}{white}{udcgray!25}
\small
\setlength{\tabcolsep}{3pt}
\begin{longtable}{ p{2cm}| p{3cm}| p{2cm}| p{5.5cm} }
  \caption{Casos de uso de información meteorológica y registros}
  \label{tab:cu-informacion-meteorologia-registros} \\

  \rowcolor{udcpink!25}
  \textbf{Identificador} & \textbf{Caso de uso} & \textbf{Estado} & \textbf{Descripción} \\\hline
  \endfirsthead
  \multicolumn{4}{c}{\tablename\ \thetable{} -- {\small \textit{(viene de la página anterior)}}} \\
  \rowcolor{udcpink!25}
  \textbf{Identificador} & \textbf{Caso de uso} & \textbf{Estado} & \textbf{Descripción} \\\hline
  \endhead
  \multicolumn{4}{c}{\dotfill{\small \textit{(continúa en la página siguiente)}}\dotfill} \\
  \endfoot
  \endlastfoot

  CU-33 & Visualizar información meteorológica relativa a una emergencia & Actualizado & Un miembro de equipo, un responsable de equipo o un coordinador podrá consultar la información meteorológica asociada a una emergencia activa. Este caso se ha ampliado para incorporar el contexto operativo de la nueva plataforma. \\
  CU-34 & Visualizar históricos de una emergencia extinguida & Actualizado & Un miembro de equipo, un responsable de equipo o un coordinador podrá consultar el histórico de una emergencia ya resuelta. Este caso se ha ampliado para mejorar la trazabilidad del incidente. \\
  CU-35 & Obtener estadísticas globales de una emergencia extinguida & Actualizado & Un miembro de equipo, un responsable de equipo o un coordinador podrá consultar las estadísticas globales de una emergencia ya resuelta. Este caso se ha ampliado para ofrecer una visión más completa del cierre operativo. \\
\end{longtable}
